% from ega3-1-fr.tex
\oldpage[0\textsubscript{III}]{21}
polynômes ayant leurs coefficients dans l'idéal $\mathfrak J$~; il est clair
que $A'/\mathfrak J'$ est canoniquement isomorphe à $k[T]$. Montrons que
l'anneau de fractions $B=A'_{\mathfrak J'}$ répond à la question~; c'est
évidemment un anneau local noethérien dont l'idéal maximal
$\mathfrak L=\mathfrak J B$. En outre,
$B/\mathfrak L=(A'/\mathfrak J')_{\mathfrak J'}=(k[T])_{\mathfrak J'}$
n'est autre que le corps des fractions $K$ de $k[T]$. Enfin, $B$ est un
$A'$-module plat et $A'$ un $A$-module libre, donc $B$ est un $A$-module plat
(\hyperref[0.6.2.1-fr]{$0_{\mathrm I}$, 6.2.1}).
\end{env}

\begin{env}[10.3.1.2]
\label{0.10.3.1.2-fr}
Supposons ensuite que $K=k(t)=k[t]$, où $t$ est algébrique sur $k$~; soit
$f\in k[T]$ le polynôme minimal de $t$~; il existe un polynôme unitaire
$F\in A[T]$ dont l'image canonique dans $k[T]$ soit $f$. Posons encore
$A'=A[T]$, et soit $\mathfrak J'$ l'idéal $\mathfrak J A'+(F)$ dans $A'$.
Nous allons voir que l'anneau quotient $B=A'/(F)$ répond cette fois à la
question. Tout d'abord, il est clair que $B$ est un $A$-module \emph{libre},
donc plat. L'anneau $A'/\mathfrak J'$ est isomorphe à
\[
  (A'/\mathfrak J A')/((\mathfrak J A'+(F))/\mathfrak J A')
  =k[T]/(f)=K~;
\]
l'image $\mathfrak L$ de $\mathfrak J'$ dans $B$ est donc maximal et on a
évidemment $\mathfrak L=\mathfrak J B$. Enfin, $B$ est un anneau semi-local,
puisqu'il est un $A$-module de type fini (Bourbaki, \emph{Alg. comm.}, chap.
IV, \S~2, n\textsuperscript{o}~5, cor. 3 de la prop. 9), et ses idéaux
maximaux sont en correspondance biunivoque avec ceux de $B/\mathfrak J B$
([13], vol. I, p. 259)~; ce qui précède prouve donc que $B$ est un anneau
local.
\end{env}

\begin{lemma}[10.3.1.3]
\label{0.10.3.1.3-fr}
Soit $(A_\lambda,f_{\mu\lambda})$ un système inductif filtrant d'anneaux
locaux, tel que les $f_{\mu\lambda}$ soient des homomorphismes locaux~; soit
$\mathfrak m_\lambda$ l'idéal maximal de $A_\lambda$, et soit
$K_\lambda=A_\lambda/\mathfrak m_\lambda$. Alors
$A'=\varinjlim A_\lambda$ est un anneau local dont
$\mathfrak m'=\varinjlim\mathfrak m_\lambda$ est l'idéal maximal, et
$K=\varinjlim K_\lambda$ le corps résiduel. En outre, si
$\mathfrak m_\mu=\mathfrak m_\lambda A_\mu$ pour $\lambda<\mu$, on a
$\mathfrak m'=\mathfrak m_\lambda A'$ pour tout $\lambda$. Si, de plus, pour
$\lambda<\mu$, $A_\mu$ est un $A_\lambda$-module plat, et si tous les
$A_\lambda$ sont noethériens, alors $A'$ est noethérien et est un
$A_\lambda$-module plat pour tout $\lambda$.
\end{lemma}

Comme $f_{\mu\lambda}(\mathfrak m_\lambda)\subset\mathfrak m_\mu$ pour
$\lambda<\mu$ par hypothèse, les $\mathfrak m_\lambda$ forment un système
inductif, et sa limite $\mathfrak m'$ est évidemment un idéal de $A'$. En
outre, si $x'\notin\mathfrak m'$, il existe $\lambda$ tel que
$x'=f_\lambda(x_\lambda)$ pour un $x_\lambda\in A_\lambda$
($f_\lambda:A_\lambda\to A'$ désignant l'homomorphisme canonique)~; puisque
$x'\notin\mathfrak m'$, on a nécessairement
$x_\lambda\notin\mathfrak m_\lambda$, donc $x_\lambda$ admet un inverse
$y_\lambda$ dans $A_\lambda$, et $y'=f_\lambda(y_\lambda)$ est l'inverse de
$x'$ dans $A'$, ce qui prouve que $A'$ est un anneau local d'idéal maximal
$\mathfrak m'$~; l'assertion relative à $K$ résulte aussitôt du fait que
$\varinjlim$ est un foncteur exact. L'hypothèse
$\mathfrak m_\mu=\mathfrak m_\lambda A_\mu$ signifie que l'application
canonique
$\mathfrak m_\lambda\otimes_{A_\lambda}A_\mu\to\mathfrak m_\mu$ est
surjective~; la relation $\mathfrak m'=\mathfrak m_\lambda A'$ résulte donc
encore de l'exactitude du foncteur $\varinjlim$ et du fait qu'il commute avec
le produit tensoriel.

Supposons maintenant que pour $\lambda<\mu$, on ait
$\mathfrak m_\mu=\mathfrak m_\lambda A_\mu$ et que $A_\mu$ soit un
$A_\lambda$-module plat. Alors $A'$ est un $A_\lambda$-module plat pour tout
$\lambda$, en vertu de (\hyperref[0.6.2.3-fr]{$0_{\mathrm I}$, 6.2.3})~;
comme $A'$ et $A_\lambda$ sont des anneaux locaux et que
$\mathfrak m'=\mathfrak m_\lambda A'$, $A'$ est même un
$A_\lambda$-module \emph{fidèlement plat}
(\hyperref[0.6.6.2-fr]{$0_{\mathrm I}$, 6.6.2}). Supposons enfin, en outre,
les $A_\lambda$ \emph{noethériens}~; les topologies
$\mathfrak m_\lambda$-préadiques sont alors séparées
(\hyperref[0.7.3.5-fr]{$0_{\mathrm I}$, 7.3.5})~; montrons qu'il en résulte
d'abord que sur $A'$ la topologie $\mathfrak m'$-adique est \emph{séparée}.
En effet, si $x'\in A'$ appartient à tous les $\mathfrak m'^{n}$ $(n>0)$,
il est l'image d'un $x_\mu\in A_\mu$ pour un certain indice $\mu$, et comme
l'image réciproque dans $A_\mu$ de
$\mathfrak m'^{n}=\mathfrak m_\mu^n A'$ est $\mathfrak m_\mu^n$
(\hyperref[0.6.6.1-fr]{$0_{\mathrm I}$, 6.6.1}), $x_\mu$ appartient à tous
les $\mathfrak m_\mu^n$, donc $x_\mu=0$ par hypothèse, et par conséquent
$x'=0$. Notons $\widehat A'$ le complété de $A'$ pour la topologie
$\mathfrak m'$-adique~; ce qui précède montre que l'on a
$A'\subset\widehat A'$. Nous allons montrer que $\widehat A'$ est
\emph{noethérien} et $A_\lambda$-\emph{plat} pour tout $\lambda$~; il en
